% PURPOSE: State formal propositions implied by ETS.K_EA.v1.1 phase logic.
% TARGET: adversarial referees.
% LANGUAGE MODE: formal.
% EVIDENCE PROFILE: internal logical consistency only (ETS-faithful).
% ROLE: claims layer (diagnostic, non-interventional).

\section{Propositions on threshold-induced transitions}
\label{sec:props}

\begin{proposition}[Uniqueness of phase classification]
\label{prop:unique-phase}
For every admissible state $s \in \Omega(K_{EA})$, the phase classification
\[
\Pi(s) \in \{\texttt{SUCCESS}, \texttt{INERTIA}, \texttt{COLLAPSE}, \texttt{STOP}\}
\]
is uniquely determined by the ETS phase predicates together with the fixed
precedence ordering.
\end{proposition}

\begin{proof}[Sketch]
By Definition~\ref{sec:phases}, each phase predicate is Boolean-valued.
If more than one predicate evaluates to true at $s$, the total precedence
ordering
\[
\texttt{STOP} \succ \texttt{COLLAPSE} \succ \texttt{INERTIA} \succ \texttt{SUCCESS}
\]
selects exactly one phase label. Hence $\Pi(s)$ is well-defined and unique.
\end{proof}

\begin{proposition}[Discontinuous reassignment under threshold crossings]
\label{prop:discontinuous}
Let $s(\lambda)$ be a parametrized family of states in $\Omega(K_{EA})$.
If, for some $\lambda^\ast$, the truth value of an ETS phase predicate
changes due to a threshold crossing (e.g.\ $f_k(s(\lambda))$ changing sign
or loss of observability), then the induced phase classification
$\Pi(s(\lambda))$ may change discontinuously at $\lambda^\ast$.
\end{proposition}

\begin{proof}[Sketch]
Phase predicates are defined by strict Boolean conditions on ETS variables.
A change in predicate truth value corresponds to crossing the boundary between
disjoint subsets of $\Omega(K_{EA})$. No continuity requirement on the parameter
$\lambda$ is imposed by ETS.
\end{proof}

\begin{proposition}[STOP as an absorbing phase]
\label{prop:stop-absorbing}
If $\Pi(s)=\texttt{STOP}$ for some $s \in \Omega(K_{EA})$, then no ETS-allowed
phase transition exists from $s$ to any other phase class.
\end{proposition}

\begin{proof}[Sketch]
The ETS allowance relation specifies $\texttt{STOP} \Rightarrow \varnothing$.
Therefore, no outgoing phase-level transitions are admissible from STOP.
\end{proof}

\begin{proposition}[Irreversibility of COLLAPSE at the phase level]
\label{prop:collapse-no-return}
If $\Pi(s)=\texttt{COLLAPSE}$, then no ETS-allowed phase transition leads
directly from COLLAPSE to either SUCCESS or INERTIA.
\end{proposition}

\begin{proof}[Sketch]
According to ETS, the only admissible outgoing transition from COLLAPSE is
$\texttt{COLLAPSE} \Rightarrow \{\texttt{STOP}\}$. No other phase-level
transitions are permitted.
\end{proof}

\begin{proposition}[Diagnostic non-steerability]
\label{prop:non-steerable}
Phase classification under ETS does not imply the existence of a continuous,
controllable, or reversible path between phase classes within
$\Omega(K_{EA})$.
\end{proposition}

\begin{proof}[Sketch]
ETS defines phase transitions as a guarded allowance relation, not as a
dynamical or control system. Consequently, neither reachability nor
steerability follows from phase identification alone.
\end{proof}

\begin{remark}[Non-claims]
The propositions above do not introduce optimization objectives, recovery
strategies, governance prescriptions, or continuity assumptions. They are
logical consequences of ETS predicates, precedence, and allowance rules only.
\end{remark}
